\section{Inference-time：Test-time Scaling 的工程价值}

\subsection{核心思想：把额外计算放在最需要的样本上}
讲者把测试时扩展看作第三种杠杆：不改参数，只在推理时增加采样、校验、检索或反思步骤，让难题获得更多计算预算。这种策略通常对长链推理最有效。

\begin{importantbox}{Test-time Scaling 的三类常用策略}
\begin{itemize}
    \item \textbf{Best-of-N / self-consistency}：多样采样后通过规则或模型选择；
    \item \textbf{检索增强迭代}：在推理中间引入外部证据，动态修正；
    \item \textbf{过程约束搜索}：用过程评分或约束过滤中间轨迹。
\end{itemize}
\end{importantbox}

\begin{figure}[H]
\centering
\includegraphics[width=0.9\textwidth]{frame-07-testtime-scaling.jpg}
\caption{视频约 01:12:40：测试时扩展在推理任务上的收益与成本权衡}
\end{figure}

\subsection{部署侧权衡：准确率、延迟与成本}
\begin{table}[H]
\centering
\renewcommand{\arraystretch}{1.2}
\begin{tabular}{p{3.0cm}p{3.4cm}p{3.4cm}p{2.8cm}}
\toprule
\textbf{策略} & \textbf{准确率变化} & \textbf{延迟/成本变化} & \textbf{适合场景} \\
\midrule
单次解码 & 基线 & 最低 & 在线高并发场景 \\
Best-of-N & 通常明显提升 & 随 N 线性增加 & 高价值低频请求 \\
自一致性 + 验证 & 对可验证任务收益大 & 中高 & 数学/代码/结构化推理 \\
检索迭代（Self-RAG） & 事实性与可追溯性提升 & 依赖检索与重排成本 & 知识密集问答 \\
\bottomrule
\end{tabular}
\caption{测试时扩展的部署权衡}
\end{table}

\begin{warningbox}{测试时扩展不是免费午餐}
若基础模型候选质量差，多采样只会放大错误；若验证器不可靠，选择过程会将噪声当信号。因此 test-time scaling 的上限受底座质量与校验机制双重约束。
\end{warningbox}

\begin{figure}[H]
\centering
\includegraphics[width=0.9\textwidth]{frame-08-selfrag.jpg}
\caption{视频约 01:17:20：Self-RAG 风格方法体现 ``推理-检索-校验'' 的闭环思想}
\end{figure}

\subsection{本章小结}
测试时扩展是部署侧的精细化预算分配工具，适合高难度、高价值请求，但必须与可靠验证机制配套使用。


